% main.tex — Gradient Inversion Attack Practice Problem
\documentclass[12pt]{article}

\usepackage[utf8]{inputenc}
\usepackage[T1]{fontenc}
\usepackage{lmodern}
\usepackage[english]{babel}
\usepackage{geometry}
\geometry{a4paper, margin=2.5cm}

\usepackage{amsmath, amssymb, amsfonts}
\usepackage{enumitem}
\usepackage{xcolor}
\usepackage{hyperref}
\usepackage{listings}
\lstset{
  basicstyle=\ttfamily\small,
  frame=single,
  breaklines=true,
  tabsize=2
}

\title{Practice Problem Part I — Non-Linear Optimization\\Gradient Inversion Attack in Federated Learning}
\author{Non-Linear Optimization Course}
\date{\today}

\begin{document}
\maketitle
\section*{Background}
Federated Learning (FL) keeps training data on client devices and shares only gradients
\(\mathbf g = \nabla_{\boldsymbol\theta}\mathcal L\bigl(f_{\boldsymbol\theta}(\mathbf x),y\bigr)\).
Yet, \emph{gradient inversion} can reconstruct \((\mathbf x,y)\) from \(\mathbf g\) and the current model
\(f_{\boldsymbol\theta}\)~\cite{zhu2019dlg,zhu2020idlg,zhu2024dlgfb,guo2025survey}.

\section*{Optimization Problem}
Given:
\begin{itemize}[nosep]
  \item model parameters \(\boldsymbol\theta\) (2-layer MLP, trained on MNIST);
  \item a leaked gradient vector \(\mathbf g^\star\) from a \emph{single} unknown sample.
\end{itemize}
Recover an image \(\hat{\mathbf x}\!\in[0,1]^{28\times28}\) and label
\(\hat y\!\in\!\{0,\dots,9\}\) by minimizing
\[
\min_{\mathbf x\in[0,1]^{784},\,y\in\{0,\dots,9\}}
\Bigl\lVert
\nabla_{\boldsymbol\theta}\mathcal L\bigl(f_{\boldsymbol\theta}(\mathbf x),y\bigr)
-\mathbf g^\star
\Bigr\rVert_2^2
+\lambda\|\mathbf x\|_2^2 .
\]

\section*{Tasks}
\begin{enumerate}[label=\textbf{\arabic*.}, nosep]
  \item Implement L-BFGS, Adam, or both, to optimize \(\mathbf x\) (continuous) and \(y\)
        (use one-hot logits or the iDLG pseudo-label trick).
  \item Run at least five random initializations and keep the best solution.
  \item Plot objective value over iterations and pixel MSE versus ground truth
        (provided for grading).
  \item Study sensitivity to \(\lambda\) and optimizer hyper-parameters.
\end{enumerate}

\section*{Deliverables (100 pts)}
\begin{itemize}[nosep]
  \item \textbf{Code (40 pts)} — reproducible, with \texttt{requirements.txt}.
  \item \textbf{Report (40 pts)} — max 4 pages: formulation, method, results, discussion.
  \item \textbf{Slides (20 pts)} — 5 min lightning talk.
\end{itemize}

\section*{Starter Code (PyTorch)}
\begin{lstlisting}[language=Python]
import torch, torch.nn as nn, torch.optim as optim

# resources provided
theta = torch.load("mlp_mnist_weights.pt")
g_star = torch.load("one_sample_gradient.pt")

model = nn.Sequential(
    nn.Flatten(),
    nn.Linear(784, 128), nn.ReLU(),
    nn.Linear(128, 10)
)
with torch.no_grad():
    for p, src in zip(model.parameters(), theta):
        p.copy_(src)

loss_fn = nn.CrossEntropyLoss()

def reconstruct(lambda_reg=1e-4, iters=300, lr=0.1, seed=0):
    torch.manual_seed(seed)
    x_hat = torch.rand(1, 1, 28, 28, requires_grad=True)
    y_logits = torch.zeros(1, 10, requires_grad=True)
    opt = optim.LBFGS([x_hat, y_logits], lr=lr, max_iter=20)

    for _ in range(iters):
        def closure():
            opt.zero_grad()
            y_prob = y_logits.softmax(dim=-1)
            loss = loss_fn(model(x_hat), y_prob)
            grads = torch.autograd.grad(loss, model.parameters(),
                                        create_graph=True)
            grad_diff = torch.cat([g.view(-1) for g in grads]) - g_star
            obj = grad_diff.pow(2).sum() + lambda_reg * x_hat.pow(2).sum()
            obj.backward()
            return obj
        opt.step(closure)
    return x_hat.detach(), y_logits.argmax(dim=-1).item()
\end{lstlisting}

\section*{Hints}
\begin{itemize}[nosep]
  \item Normalise \(\mathbf g^\star\) and use gradient clipping.
  \item Use multiple restarts—the landscape has many poor local minima.
  \item TV-loss or generative priors (GAN latents) can refine images but enlarge the search space.
\end{itemize}

\section*{Ethics Notice}
Apply this attack \textbf{only} to the classroom model and data.

\begin{thebibliography}{99}\small
\bibitem{zhu2019dlg} L. Zhu, Z. Liu, S. Han. \emph{Deep Leakage from Gradients}. NeurIPS 2019.
\bibitem{zhu2020idlg} L. Zhu, S. Han. \emph{iDLG: Improved Deep Leakage from Gradients}. arXiv:2001.02610 (2020).
\bibitem{zhu2024dlgfb} A. Ribeiro \emph{et al.} \emph{DLG-FB: Feedback-Blending Gradient Inversion}. arXiv:2409.17767 (2024).
\bibitem{guo2025survey} P. Guo \emph{et al.} \emph{Survey of Gradient Inversion Attacks}. 2025.
\end{thebibliography}

\end{document}
